%%%%%%%%%%%%%%%%%%%%%%%%%%%%%%%%%%%%%%%%%%%%%%%%%%%%%%%%%%%%%%%%%%%%%%%%%%%
%%  Project: CPHOS LaTeX Template
%%  File: example-paper.tex
%%  Copyright © 2026 gameswu & CPHOS (www.cphos.cn)
%%
%%  Permission is hereby granted, free of charge, to any person obtaining a 
%%  copy of this software and associated documentation files (the "Software"), 
%%  to deal in the Software without restriction, including without limitation 
%%  the rights to use, copy, modify, merge, publish, distribute, sublicense, 
%%  and/or sell copies of the Software, and to permit persons to whom the 
%%  Software is furnished to do so, subject to the following conditions:
%%
%%  The above copyright notice and this permission notice shall be included 
%%  in all copies or substantial portions of the Software.
%%
%%  THE SOFTWARE IS PROVIDED "AS IS", WITHOUT WARRANTY OF ANY KIND, EXPRESS 
%%  OR IMPLIED, INCLUDING BUT NOT LIMITED TO THE WARRANTIES OF MERCHANTABILITY, 
%%  FITNESS FOR A PARTICULAR PURPOSE AND NONINFRINGEMENT. IN NO EVENT SHALL 
%%  THE AUTHORS OR COPYRIGHT HOLDERS BE LIABLE FOR ANY CLAIM, DAMAGES OR OTHER 
%%  LIABILITY, WHETHER IN AN ACTION OF CONTRACT, TORT OR OTHERWISE, ARISING 
%%  FROM, OUT OF OR IN CONNECTION WITH THE SOFTWARE OR THE USE OR OTHER 
%%  DEALINGS IN THE SOFTWARE.
%%%%%%%%%%%%%%%%%%%%%%%%%%%%%%%%%%%%%%%%%%%%%%%%%%%%%%%%%%%%%%%%%%%%%%%%%%%

% CPHOS理论试题模板

\documentclass[exam]{cphos}

% === 设置文档信息 ===
\cphostitle{第XX届CPHOS物理竞赛联考}
\cphossubtitle{理论试题参考答案及评分标准}
\cphoswatermark{../assets/watermark.jpg}

% === 启用分值检查（可选） ===
\setscorecheck{true}  % 启用后，如果声明分值与实际分值不符会显示红色警告

\begin{document}

% === 标题页 ===
\maketitlepage

{\kaiti\small 本文件于YYYY年MM月DD日HH:MM发布，最后更新于\compiletime 。}

{\kaiti\small CPHOS物理竞赛联考是开放性公益性的考试，有意向参与的教师和学生可以关注“CPHOS”微信公众号进行报名，报名后方可参与联考。请使用“CPHOS物理竞赛联考”微信小程序完成答题卡上传、阅卷、成绩查询等操作。联系方式见试题末尾。}
% === 时间安排（仅exam模式显示）===
\begin{examschedule}
\begin{tabular}{@{}cc@{}}
    答题卡上传 & 阅卷 \\
    YYYY/MM/DD HH:MM - YYYY/MM/DD HH:MM & YYYY/MM/DD HH:MM - YYYY/MM/DD HH:MM \\
\end{tabular}

\vspace{0.3em}

\begin{tabular}{@{}ccc@{}}
    非正式成绩 & 成绩申诉 & 正式成绩 \\
    YYYY/MM/DD HH:MM & YYYY/MM/DD HH:MM - YYYY/MM/DD HH:MM & YYYY/MM/DD HH:MM \\
\end{tabular}
\end{examschedule}

% === 考生须知（仅exam模式显示）===
\begin{examnotes}
    \item 理论试题共 \textbf{\underline{\cphostotalpages}} 页，理论答题卡共 \textbf{\underline{X}} 页，答题时间 \textbf{\underline{180}} 分钟，试题满分 \textbf{\underline{\cphostotalscore}} 分。
    \item 请在答题卡的指定答题区域内答题，试题和草稿纸上的内容将不会作为评分参考，不可申请答题卡加页。
    \item 若发现试题存在问题，请向领队（教练）反映，由其转达至相关微信群聊。
    \item 试题答案及相关分析均会在官方网站 \url{www.cphos.cn} 上发布。
    \item 本次考试定位难度为决赛。
\end{examnotes}
\vspace{1em}

% ============================================
% 第一题
% ============================================
\begin{problem}[40]{太阳物理初步}

\begin{problemstatement}
太阳与我们的生活息息相关，本题将逐步深入来研究太阳的一些基本物理性质，由于太阳非常复杂所以在很多地方我们将做一些非常粗糙的近似，不过这不会影响我们最终结果的正确性。本题中可能会用到的一些物理常数有，日地距离$r_{es}=1.50\times 10^{11}\text{~m}$，太阳半径$R_s=6.96\times 10^{8}\text{~m}$，万有引力常数$G=6.67\times 10^{-11}\text{~m}^3\cdot\text{kg}^{-1}\cdot\text{s}^{-2}$，太阳光度即太阳单位时间对外辐射的能量$L_s=3.90\times 10^{26}\text{~J}\cdot\text{s}^{-1}$，地球绕太阳的公转周期$t=365\text{~d}$，斯特藩常数$\sigma=5.67\times 10^{-8}\text{~W}\cdot\text{m}^{-2}\cdot\text{K}^{-4}$，韦恩位移常数$b=2.90\times 10^{-3}\text{~m}\cdot\text{K}$，玻尔半径$r_1=5.29\times 10^{-11}\text{~m}$。

\subq{1}利用题目给出的物理常数计算太阳的质量$M_s$，平均密度$\overline{\rho}_s$以及太阳的表面温度$T_s$。同时作为一种近似（尽管事实并非如此），假定太阳的密度处处相同并且为上一小问得到的平均密度$\overline{\rho}_s$，同时由于太阳整体没有明显的膨胀和收缩现象，因此可以使用流体静力学方程来处理，据此给出太阳内部的压强分布$P(r)$。

\subq{2}在本问中我们不再认为太阳是一个密度均匀的球体，并将考虑太阳内部的一些物理性质。

\subsubq{2.1}作为一种近似模型，我们假定太阳内部元素只有氢原子（未电离），并且氢原子按照立方体晶格的形式进行排列，每个氢原子排列在边长为$2r_1$（$r_1$是题目中给出的玻尔半径）的立方体的一个顶点上，见图1.1，据此求出太阳内部的平均密度$\rho_s'$并将其与第（1）问中求出的平均密度进行比较。

\subsubq{2.2}实际计算表明太阳中心的密度约为$\rho_c=1.5\times 10^5\text{~kg}\cdot\text{m}^{-3}$，据此讨论太阳中心的物质状态。
\end{problemstatement}

\begin{solution}

% === 使用\solsubq自动记分 ===
\solsubq{1}{10}
考虑地球绕太阳进行公转，由牛顿第二定律可以得到
\begin{equation}
    \frac{GM_s m}{r_{es}^2} = m \frac{4\uppi^2}{t^2} r_{es} \eqtagscore{1}{2}
\end{equation}
由此解出
\begin{equation}
    M_s = \frac{4\uppi^2 r_{es}^3}{Gt^2} = 2.01 \times 10^{30} \text{~kg} \eqtagscore{2}{2}
\end{equation}
\begin{equation}
    \overline{\rho}_s = \frac{M_s}{\frac{4}{3}\uppi R_s^3} = 1.42 \text{~g}\cdot\text{cm}^{-3} \eqtagscore{3}{2}
\end{equation}
由斯特藩定律也可以直接得到
\begin{equation}
    T_s = \sqrt[4]{\frac{L_s}{4\uppi\sigma R_s^2}} = 5.80 \times 10^{3} \text{~K} \eqtagscore{4}{2}
\end{equation}
而对于位于$r\sim r+\mathrm{d}r$的微元来说，径向的压强差正好和其引力相平衡，由此得到
\begin{equation}
    \mathrm{d}P(r) = \frac{GM(r)\rho_s}{r^2}\,\mathrm{d}r = \frac{4}{3}\uppi G\rho_s^2 r\,\mathrm{d}r \eqtag{5}
\end{equation}
对于太阳表面来说，由于外部是极其稀薄的太空，因此压强近似为0，由此直接积分可以得到
\begin{equation}
    P(r) = \frac{2}{3}\uppi G\rho_s^2(R_s^2 - r^2) \eqtagscore{6}{2}
\end{equation}

\solsubq{2}{15}
\solsubsubq{2.1}{3}
根据题意，每个氢原子占据边长为$2r_1$的立方体，则
\begin{equation}
    \rho_s' = \frac{m_H}{(2r_1)^3} \eqtagscore{7}{2}
\end{equation}
代入数据可得$\rho_s' \approx 2.8\text{~g}\cdot\text{cm}^{-3}$，与第（1）问结果同一数量级。\addtext{讨论}{1}

\solsubsubq{2.2}{6}
根据题目给出的太阳中心密度，可以计算出
\begin{equation}
    n = \frac{\rho_c}{m_H} \eqtagscore{8}{2}
\end{equation}
由此得到原子间距
\begin{equation}
    d = n^{-1/3} \approx 2.5 \times 10^{-11} \text{~m} \eqtagscore{9}{2}
\end{equation}
这个距离小于玻尔半径，说明太阳中心的物质处于等离子体状态。\addtext{讨论}{2}

% === 自动生成评分标准 ===
\scoring

\end{solution}

% === 题目附录（可选） ===
\begin{problemappendix}
本题考查了恒星物理的基本概念，包括光度、有效温度、流体静力学平衡等。学生需要熟悉斯特藩-玻尔兹曼定律、理想气体状态方程等基本物理关系。
\end{problemappendix}

\end{problem}

% ============================================
% 版权信息（使用copyrightinfo环境确保不分页）
% ============================================
\begin{copyrightinfo}
\begin{center}
    {\textbf{版权信息}}
\end{center}

\noindent{\textbf{命题人}}

X~X\quad XXX\quad XXX

\noindent{\textbf{审题人}}

Y~Y\quad YYY\quad YYY

\vspace{0.5em}
\contactinfo{../assets/fig1.jpg}{../assets/fig2.png}{../assets/fig3.png}{\url{service@cphos.cn}}{CPHOS物理竞赛联考}
\end{copyrightinfo}

\end{document}
